\documentclass[11pt]{article}
\usepackage[letterpaper,margin=1in]{geometry}
\usepackage[T1]{fontenc}
\usepackage[utf8]{inputenc}
\usepackage{lmodern,microtype}
\usepackage{amsmath,amssymb,mathtools,bm,slashed}
\usepackage{xcolor,enumitem,booktabs,array,needspace}
\usepackage[most]{tcolorbox}
\usetikzlibrary{decorations.pathmorphing,arrows.meta}
\tikzset{electron/.style={thick,-{Stealth[length=2mm]}},
 photon/.style={thick,decorate,decoration={snake,amplitude=1pt,segment length=4pt}}}
\usepackage{fancyhdr,titlesec,parskip,hyperref}
\definecolor{qftblue}{HTML}{245A8D}
\definecolor{qftgreen}{HTML}{2F6B4F}
\definecolor{qftgray}{HTML}{F4F6F8}
\hypersetup{colorlinks=true,linkcolor=qftblue,urlcolor=qftblue,pdfauthor={Nima Lashkari}}
\pagestyle{fancy}\fancyhf{}
\lhead{PHYS 662: Quantum Field Theory I}\cfoot{\thepage}
\setlength{\headheight}{14pt}
\titleformat{\section}{\Large\bfseries\color{qftblue}}{\thesection.}{0.6em}{}
\titleformat{\subsection}{\large\bfseries\color{qftgreen}}{\thesubsection}{0.6em}{}
\setlist[enumerate]{leftmargin=2em,itemsep=0.25em,topsep=0.35em}
\setlist[itemize]{leftmargin=2em,itemsep=0.2em,topsep=0.35em}
\newcommand{\dd}{\mathrm d}
\newcommand{\ii}{\mathrm i}
\newcommand{\eps}{\varepsilon}
\newcommand{\tr}{\operatorname{tr}}
\newcommand{\Tr}{\operatorname{Tr}}
\newcommand{\diag}{\operatorname{diag}}
\newcommand{\Lag}{\mathcal L}
\newcommand{\D}{\mathcal D}
\newcommand{\vac}{\lvert0\rangle}
\newcommand{\vbra}{\langle0\rvert}
\newcommand{\Tord}{\mathrm T}
\newcommand{\ket}[1]{\lvert#1\rangle}
\newcommand{\bra}[1]{\langle#1\rvert}
\newtcolorbox{keyidea}[1][]{enhanced,breakable,colback=qftgray,colframe=qftblue,
 boxrule=0.7pt,arc=1.5mm,title={#1},fonttitle=\bfseries,coltitle=white}
\newtcolorbox{remarkbox}[1][]{enhanced,breakable,colback=white,colframe=qftgreen,
 boxrule=0.6pt,arc=1.5mm,title={#1},fonttitle=\bfseries,coltitle=white}
\newcommand{\intp}{\int\!\frac{\dd^4p}{(2\pi)^4}}
\newcommand{\MSbar}{\overline{\mathrm{MS}}}
\newcommand{\module}[1]{\clearpage\section{#1}}
\newcommand{\braket}[2]{\langle#1\mid#2\rangle}
\newcommand{\dPi}{\dd\Pi}
\newcommand{\cM}{\mathcal M}
\newcommand{\cH}{\mathcal H}
\newcommand{\cS}{\mathcal S}
\newcommand{\Zp}{Z_{\rm pole}}
\newcommand{\Kallen}{\mathcal K}
\newcommand{\lecturefront}[4]{%
\rhead{Lecture #1}\hypersetup{pdftitle={PHYS 662 Lecture #1: #2}}
\begin{document}\pagenumbering{gobble}
\begin{titlepage}\centering\vspace*{1.05in}
{\scshape\Large PHYS 662 --- Quantum Field Theory I\par}
\vspace{0.35in}{\Huge\bfseries\color{qftblue}Lecture #1\par}
\vspace{0.2in}{\LARGE\bfseries #2\par}
\vspace{0.2in}{\large\color{qftgreen}#3\par}
\vspace{0.6in}{\large Nima Lashkari\par}{\large Fall 2026\par}
\vfill\begin{minipage}{0.86\textwidth}\small #4\par\medskip
Course page: \url{https://nimalashkari.com/quantum-field-theory-2026-part-i/}
\end{minipage}\vspace{0.5in}\end{titlepage}
\pagenumbering{roman}\setcounter{tocdepth}{1}\tableofcontents
\vfill\begin{remarkbox}[Place in the course]
Lecture 26 constructs the physical photon Hilbert space and its propagators.
Lecture 27 develops QED amplitudes. Lecture 28 computes the photon and
fermion self-energies. Lecture 29 relates the vertex to charge
renormalization and the magnetic moment. Homework 9 closes this sequence.
\end{remarkbox}
\small\textbf{Reading and conventions.} The uploaded course slides are the
primary starting point. D.~Tong,
\href{https://www.damtp.cam.ac.uk/user/tong/qft/qfthtml/S6.html}{\textit{Quantum Field Theory}, Chapter 6},
is a companion for photon quantization and QED. For the vertex calculation,
see V.~Kaplunovsky,
\href{https://web2.ph.utexas.edu/~vadim/Classes/2020f/vertex.pdf}{\textit{QED Vertex Correction}},
especially the parameter integral for the Pauli form factor.
For further textbook study, see Peskin--Schroeder, \textit{An Introduction
to QFT}, Chapters 5--7 and 10.
All formulas here use the mostly-minus metric and propagators that include
\(\ii\). Dimensional regularization uses \(d=4-2\eps\), which differs from
some of the slides. The accompanying editorial memo records corrections
and the scope of the supplementary reading.
\normalsize\clearpage\pagenumbering{arabic}}
\lecturefront{27}{Quantum Electrodynamics}
{Gauge invariance, scattering amplitudes, and the organization of loops}
{We couple the quantized electromagnetic field to Dirac matter and turn
the resulting action into measurable scattering probabilities. Worked
examples make the roles of spin, identical particles, and the Ward
identity explicit, before we organize the loop corrections studied in
the final two lectures.}

\section{Minimal coupling and the conserved current}
Continue with \(\eta=\diag(1,-1,-1,-1)\), and let \(e\) be the signed
charge of the Dirac particle. Thus \(e<0\) for an electron, while
\(\alpha=e^2/(4\pi)>0\). The gauge-fixed QED Lagrangian is
\begin{equation}
 \Lag=-\frac14F_{\mu\nu}F^{\mu\nu}
 -\frac{(\partial\cdot A)^2}{2\xi}
 +\bar\psi(\ii\slashed\partial-m)\psi
 -e\bar\psi\gamma^\mu\psi A_\mu.
 \label{eq:qed}
\end{equation}
Its gauge-invariant part follows by replacing \(\partial_\mu\) with
\(D_\mu=\partial_\mu+\ii e A_\mu\). Indeed,
\begin{equation}
 A_\mu\longmapsto A_\mu+\partial_\mu\chi,\qquad
 \psi\longmapsto e^{-\ii e\chi}\psi,
 \qquad D_\mu\psi\longmapsto e^{-\ii e\chi}D_\mu\psi.
\end{equation}
The term involving \(\partial_\mu\chi\) from the ordinary derivative
cancels the new gauge-potential term. This local cancellation fixes the
relative coefficient of the kinetic and interaction terms.

The global phase transformation has current
\begin{equation}
 j^\mu=e\bar\psi\gamma^\mu\psi,\qquad \partial_\mu j^\mu=0.
\end{equation}
The interacting Dirac equation and its adjoint prove the second relation;
the terms containing \(A_\mu\) cancel. The Maxwell equation is
\(\partial_\mu F^{\mu\nu}=j^\nu\) before gauge fixing. An antiparticle
carries charge \(-e\), as follows from the normal-ordered charge operator.

In four dimensions, \([A]=1\), \([\psi]=3/2\), and \([e]=0\).
The interaction is dimension four. A photon mass term would violate the
gauge redundancy, and a Pauli term \(\bar\psi\sigma^{\mu\nu}\psi F_{\mu\nu}\)
has dimension five. Neither is an independent term in minimal
renormalizable QED, although the latter will arise as a finite
low-energy description of a loop effect.

\begin{keyidea}[Symmetry relates computations]
Gauge invariance constrains amplitudes, the ultraviolet counterterms,
and charge renormalization. We will first see it as cancellation between
Compton diagrams; Lecture 29 derives the corresponding identity for
full propagators and the vertex.
\end{keyidea}
There is no elementary three- or four-photon vertex in \eqref{eq:qed}.
Photons nevertheless interact through charged-particle loops. Free
Abelian ghosts decouple, as shown in Lecture 26.

\module{Feynman rules and fermion bookkeeping}
Our convention is \(S=1+\ii T\), with connected matrix element
\(\ii(2\pi)^4\delta^{(4)}(P_f-P_i)\cM\). The following rules therefore
produce \(\ii\cM\), not \(\cM\). In Feynman gauge,
\begin{center}
\renewcommand{\arraystretch}{1.35}
\begin{tabular}{ll}
\toprule
Object & Factor \\
\midrule
Internal fermion & \(\displaystyle S_F(p)=\frac{\ii(\slashed p+m)}{p^2-m^2+\ii0}\) \\
Internal photon & \(\displaystyle D_{\mu\nu}(k)=\frac{-\ii\eta_{\mu\nu}}{k^2+\ii0}\) \\
Photon--fermion vertex & \(-\ii e\gamma^\mu\) \\
Incoming / outgoing fermion & \(u_s(p)\ /\ \bar u_s(p)\) \\
Incoming / outgoing antifermion & \(\bar v_s(p)\ /\ v_s(p)\) \\
Incoming / outgoing photon & \(\epsilon_\mu(k)\ /\ \epsilon_\mu^*(k)\) \\
\bottomrule
\end{tabular}
\end{center}
Momentum is conserved at every vertex. Follow each fermion line to form
an ordered matrix product between its external spinors. A closed line
produces a trace and an additional factor \(-1\), originating from the
cyclic rearrangement of Grassmann fields. Exchange of identical external
fermions gives a relative minus sign between the corresponding diagrams.
These are distinct bookkeeping rules.

Integrate each independent loop momentum using
\(\int\dd^4\ell/(2\pi)^4\), regulate before shifting a divergent integral,
and include the usual graph symmetry factor. Divide out vacuum bubbles.
For unpolarized probabilities, sum final spins and polarizations and
average over the initial states actually populated.

The useful completeness relations are
\begin{equation}
 \sum_su_s(p)\bar u_s(p)=\slashed p+m,\qquad
 \sum_sv_s(p)\bar v_s(p)=\slashed p-m.
\end{equation}
For a matrix \(X\), define \(\bar X=\gamma^0X^\dagger\gamma^0\). Then
\begin{equation}
 \sum_{s,s'}|\bar u_{s'}(p')Xu_s(p)|^2
 =\tr[(\slashed p'+m)X(\slashed p+m)\bar X].
 \label{eq:spintrace}
\end{equation}
This follows by multiplying the amplitude by its complex conjugate and
inserting both completeness relations. It is more reliable than expanding
individual spinors. At tree level we use four-dimensional gamma algebra;
in ultraviolet-divergent loops the dimension-dependent identities must
be retained until the pole expansion is complete.

\module{Exchange of identical fermions}
Consider M\o ller scattering,
\(e^-(p_1,s_1)e^-(p_2,s_2)\to e^-(p_3,s_3)e^-(p_4,s_4)\).
There are two ways to connect the external fermion lines. With
\(t=(p_1-p_3)^2\) and \(u=(p_1-p_4)^2\), the rules give
\begin{equation}
 \cM=e^2\left[
 \frac{[\bar u_3\gamma^\mu u_1][\bar u_4\gamma_\mu u_2]}{t}
 -\frac{[\bar u_4\gamma^\mu u_1][\bar u_3\gamma_\mu u_2]}{u}
 \right].
 \label{eq:moller}
\end{equation}
The relative minus sign enforces antisymmetry under exchange of the
complete outgoing one-particle labels. If momentum \emph{and spin} labels
coincide, the amplitude vanishes. Equal momenta alone do not require
vanishing when the spin states differ.

Each current is conserved on shell. For example,
\begin{equation}
 (p_1-p_3)_\mu\bar u_3\gamma^\mu u_1
 =\bar u_3(\slashed p_1-\slashed p_3)u_1
 =m\bar u_3u_1-m\bar u_3u_1=0.
\end{equation}
The longitudinal part of an arbitrary covariant photon propagator
therefore drops out of each exchange term separately. This particularly
simple diagram-by-diagram cancellation will not hold in every process.

For reference, a basic trace needed in spin sums is
\begin{equation}
 \tr[(\slashed a+m)\gamma^\mu(\slashed b+m)\gamma^\nu]
 =4\left[a^\mu b^\nu+a^\nu b^\mu
 -\eta^{\mu\nu}(a\cdot b-m^2)\right].
\end{equation}
It follows from the vanishing of odd traces and
\(\tr(\gamma^\alpha\gamma^\mu\gamma^\beta\gamma^\nu)
=4(\eta^{\alpha\mu}\eta^{\beta\nu}-\eta^{\alpha\beta}\eta^{\mu\nu}
+\eta^{\alpha\nu}\eta^{\mu\beta})\).
The interference of the two exchange terms is a single longer trace;
it must be kept before any spin average.

\begin{remarkbox}[Statistics and phase space are separate]
Antisymmetrizing \(\cM\) does not remove the identical-particle phase-space
factor. Integrating over all labeled final electron momenta requires
\(1/2!\). Equivalently, integrate over one nonredundant final-state region
without that factor. The initial beam flux is not divided by \(2!\).
\end{remarkbox}
At small momentum transfer, the \(t\)-channel singularity reflects the
long-range Coulomb interaction. It is not a massive resonance; finite
angular acceptance matters when quoting an integrated cross section.

\module{Compton scattering and a diagrammatic Ward identity}
For \(e^-(p)+\gamma(k)\to e^-(p')+\gamma(k')\), let
\(p+k=p'+k'\), \(s=(p+k)^2\), and \(u=(p-k')^2\).
There are two orderings of the photon vertices along the electron line:
\begin{center}
\begin{tikzpicture}[scale=.78,font=\small]
\foreach \shift/\channel/\internal/\leftph/\rightph in
 {0/s/{p+k}/{k\ \mathrm{in}}/{k'\ \mathrm{out}},
  7/u/{p-k'}/{k'\ \mathrm{out}}/{k\ \mathrm{in}}}{
 \begin{scope}[xshift=\shift cm]
 \draw[electron](-2.3,0)--(-.65,0) node[midway,below]{$p$};
 \draw[electron](-.65,0)--(.65,0) node[midway,below]{$\internal$};
 \draw[electron](.65,0)--(2.3,0) node[midway,below]{$p'$};
 \draw[photon](-1.35,1.15)--(-.65,0);
 \draw[photon](.65,0)--(1.35,1.15);
 \node[above] at(-1.35,1.15){$\leftph$};
 \node[above] at(1.35,1.15){$\rightph$};
 \fill(-.65,0)circle(2pt);\fill(.65,0)circle(2pt);
 \node at(0,-.8){$\channel$ channel};
 \end{scope}}
\end{tikzpicture}
\end{center}
\begin{align}
 \cM=-e^2\bar u(p')\left[
 \slashed\epsilon'^{\,*}\frac{\slashed p+\slashed k+m}{s-m^2}
 \slashed\epsilon
 +\slashed\epsilon\frac{\slashed p-\slashed k'+m}{u-m^2}
 \slashed\epsilon'^{\,*}\right]u(p).
 \label{eq:compton}
\end{align}
Here \(\slashed\epsilon'^{\,*}\) means the slash of the complex-conjugate
outgoing polarization. The diagrams have a relative plus sign: no
identical external fermions have been exchanged. Minimal Dirac QED has
no additional two-photon contact vertex.

The gauge test is to replace the incoming polarization by \(k\).
Set \(Q=p+k\), \(R=p-k'\), and use the stripped free propagator
\(\mathcal S_0(r)=(\slashed r-m)^{-1}\). The first term becomes
\begin{equation}
 \bar u'\slashed\epsilon'^{\,*}\mathcal S_0(Q)\slashed k u
 =\bar u'\slashed\epsilon'^{\,*}u,
 \qquad \slashed k u=(\slashed Q-m)u.
\end{equation}
For the second term, momentum conservation gives \(k=p'-R\), so
\begin{equation}
 \bar u'\slashed k\mathcal S_0(R)\slashed\epsilon'^{\,*}u
 =-\bar u'\slashed\epsilon'^{\,*}u.
\end{equation}
The two terms cancel exactly. Replacing the outgoing polarization by
\(k'\) gives the second Ward identity in the same way.

\begin{keyidea}[Test the sum before summing polarizations]
Each Compton diagram fails the external-photon Ward test. Their sum
passes. Consequently the physical polarization sum may be replaced by
\(-\eta_{\mu\nu}\) only for the complete amplitude.
\end{keyidea}
This cancellation is the elementary form of a general telescoping
identity: contracting a photon momentum into a fermion vertex gives the
difference of the inverse propagators on its two sides. Adjacent
attachments cancel, leaving endpoint terms that vanish for external
on-shell states. Off shell the endpoint terms survive and become the
Ward--Takahashi identity, rather than zero.

\module{From Compton amplitudes to the Klein--Nishina formula}
In the laboratory take \(p=(m,\bm0)\), incoming photon energy \(\omega\),
and outgoing energy \(\omega'\), with angle \(\theta\) between photons.
The condition \((p+k-k')^2=m^2\) gives
\begin{equation}
 \omega'=\frac{\omega}{1+(\omega/m)(1-\cos\theta)}.
 \label{eq:recoil}
\end{equation}
The recoil is essential except in the low-energy limit.

Write the bracket of \eqref{eq:compton} as
\(\epsilon'^*_\nu\epsilon_\mu X^{\nu\mu}\), where
\begin{equation}
 X^{\nu\mu}=\gamma^\nu\frac{\slashed p+\slashed k+m}{2p\cdot k}\gamma^\mu
 +\gamma^\mu\frac{\slashed p-\slashed k'+m}{-2p\cdot k'}\gamma^\nu.
\end{equation}
The spin and polarization average is
\begin{equation}
 \overline{|\cM|^2}=\frac{e^4}{4}
 \eta_{\mu\rho}\eta_{\nu\sigma}
 \tr[(\slashed p'+m)X^{\nu\mu}(\slashed p+m)\bar X^{\sigma\rho}].
\end{equation}
The factor \(1/4\) averages two electron spins and two photon helicities.
Expanding the two squares and their interference, reducing the gamma
traces, and using \(p\cdot k=m\omega\), \(p\cdot k'=m\omega'\), gives
\begin{equation}
 \overline{|\cM|^2}=2e^4\left[
 \frac{\omega}{\omega'}+\frac{\omega'}{\omega}-\sin^2\theta\right].
\end{equation}
For example, before inserting laboratory variables the bracket is
\(a/b+b/a+2m^2(1/a-1/b)+m^4(1/a-1/b)^2\), with
\(a=p\cdot k\), \(b=p\cdot k'\). The recoil relation turns its last two
terms into \(-\sin^2\theta\).

Integrating the two-body phase space with flux \(4m\omega\) yields
\begin{equation}
 \frac{\dd\sigma}{\dd\Omega}
 =\frac{\overline{|\cM|^2}}{64\pi^2m^2}
 \left(\frac{\omega'}{\omega}\right)^2
 =\frac{\alpha^2}{2m^2}\left(\frac{\omega'}{\omega}\right)^2
 \left[\frac{\omega}{\omega'}+\frac{\omega'}{\omega}-\sin^2\theta\right].
\end{equation}
One recoil factor arises from final-state phase space, and another from
the energy-delta-function Jacobian. The expression is the unpolarized
Klein--Nishina formula.

When \(\omega\ll m\), \(\omega'\simeq\omega\), giving the Thomson results
\begin{equation}
 \frac{\dd\sigma_T}{\dd\Omega}=\frac{\alpha^2}{2m^2}(1+\cos^2\theta),
 \qquad \sigma_T=\frac{8\pi\alpha^2}{3m^2}.
\end{equation}
The dimensions, positivity, and low-energy limit independently check the
normalization of the relativistic calculation.

\module{Annihilation into a charged fermion pair}
Consider \(e^-(p_1)e^+(p_2)\to\mu^-(p_3)\mu^+(p_4)\), with muon mass
\(M\). Neglect the incoming electron mass but retain \(M\). At tree level
there is one annihilation diagram,
\begin{equation}
 \cM=\frac{e^2}{s}
 [\bar v(p_2)\gamma^\mu u(p_1)]
 [\bar u(p_3)\gamma_\mu v(p_4)],\qquad s=(p_1+p_2)^2.
\end{equation}
The spin average factorizes into two traces:
\begin{equation}
 \overline{|\cM|^2}=\frac{e^4}{4s^2}
 \tr[\slashed p_2\gamma^\mu\slashed p_1\gamma^\nu]
 \tr[(\slashed p_3+M)\gamma_\mu(\slashed p_4-M)\gamma_\nu].
\end{equation}
The second trace has the mass term with the opposite sign to a
fermion--fermion trace because the outgoing antiparticle completeness
relation is \(\slashed p_4-M\). Contracting gives
\begin{equation}
 \overline{|\cM|^2}
 =\frac{2e^4}{s^2}\left[(M^2-t)^2+(M^2-u)^2+2M^2s\right].
\end{equation}
In the center-of-momentum frame define
\(\beta=\sqrt{1-4M^2/s}\). With \(\theta\) the angle of the outgoing
negative muon relative to the incoming electron,
\begin{equation}
 M^2-t=\frac{s}{2}(1-\beta\cos\theta),\qquad
 M^2-u=\frac{s}{2}(1+\beta\cos\theta).
\end{equation}
The two-body phase-space formula of Lecture 24 now gives
\begin{align}
 \frac{\dd\sigma}{\dd\Omega}
 &=\frac{\alpha^2\beta}{4s}
 \left[1+\cos^2\theta+(1-\beta^2)\sin^2\theta\right],\\
 \sigma(s)&=\frac{4\pi\alpha^2}{3s}\,
 \beta\left(1+\frac{2M^2}{s}\right).
 \label{eq:paircross}
\end{align}
There is no identical-final-particle divisor. At high energy the angular
distribution approaches \(1+\cos^2\theta\), and the total cross section
scales as \(1/s\). At threshold the phase-space factor \(\beta\) suppresses
the rate even though the matrix element remains finite.

\begin{remarkbox}[A bridge to the next lecture]
The factor \(\beta(1+2M^2/s)\) will reappear in the imaginary part of
photon vacuum polarization. Cutting the two fermion lines places the
same pair on shell; unitarity relates the cut loop to the probability
for producing that intermediate state.
\end{remarkbox}

\module{Why the basic loop corrections have their form}
The one-loop photon self-energy is a closed fermion loop with two photon
attachments. The electron self-energy contains one internal photon and
one internal electron. The vertex correction contains one photon and two
electron segments. All three carry two more powers of \(e\) than their
tree-level counterparts.
\begin{center}
\begin{tikzpicture}[scale=.68,font=\small]
\begin{scope}
 \draw[photon](-2,0)--(-.8,0);\draw[photon](.8,0)--(2,0);
 \draw[electron](-.8,0)arc[start angle=180,end angle=0,x radius=.8,y radius=.65];
 \draw[electron](.8,0)arc[start angle=0,end angle=-180,x radius=.8,y radius=.65];
 \node at(0,-1.15){photon self-energy};
\end{scope}
\begin{scope}[xshift=5.7cm]
 \draw[electron](-2,0)--(-.9,0);\draw[electron](-.9,0)--(.9,0);
 \draw[electron](.9,0)--(2,0);
 \draw[photon](-.9,0)to[out=70,in=110](.9,0);
 \node at(0,-1.15){electron self-energy};
\end{scope}
\begin{scope}[xshift=11.4cm]
 \draw[electron](-2,0)--(-1,0);\draw[electron](-1,0)--(0,0);
 \draw[electron](0,0)--(1,0);\draw[electron](1,0)--(2,0);
 \draw[photon](-1,0)to[out=-65,in=-115](1,0);
 \draw[photon](0,0)--(0,1.15);
 \node at(0,-1.15){vertex correction};
\end{scope}
\end{tikzpicture}
\end{center}

For a connected QED graph, let \(I_A,I_\psi,V,L\) count internal photons,
internal fermions, vertices, and loops, and let \(E_A,E_\psi\) count
external lines. At large generic loop momenta the superficial degree of
divergence is
\begin{equation}
 \mathcal D=4L-2I_A-I_\psi.
\end{equation}
Using \(L=I_A+I_\psi-V+1\),
\(2I_\psi+E_\psi=2V\), and \(2I_A+E_A=V\), obtain
\begin{equation}
 \mathcal D=4-E_A-\frac32E_\psi.
\end{equation}
Thus the photon two-point function is superficially quadratic, the
fermion two-point function linear, and the vertex logarithmic.
This is an upper bound, before symmetry constraints. The Ward identity
forces the photon self-energy to contain
\(q^2\eta^{\mu\nu}-q^\mu q^\nu\), leaving only a logarithmic scalar
coefficient and forbidding a photon-mass counterterm.

Charge conjugation supplies another constraint. Reversing a closed
fermion loop with \(n\) photon insertions changes its sign by \((-1)^n\).
After pairing the two orientations, every odd-photon loop vanishes
(Furry's theorem). The statement assumes the charge-conjugation-invariant
vacuum, with no external medium or background that breaks that symmetry.

\begin{keyidea}[Counterterm structures do not exhaust amplitudes]
The photon two-point function, fermion two-point function, and vertex
control the independent QED ultraviolet renormalizations. They do not
generate every scattering diagram. For example, light-by-light scattering
has genuine four-photon box diagrams. Their complete one-loop sum is
ultraviolet finite by the gauge identities.
\end{keyidea}
At higher orders, divergent subgraphs must also be subtracted. Power
counting is therefore a guide to possible local counterterms, not a
replacement for the subtraction of subdivergences or for Ward identities.

\module{Preparing a consistent loop expansion}
Use dimensional regularization with \(d=4-2\eps\), and write
\begin{equation}
 \frac1{\bar\eps}=\frac1\eps-\gamma_E+\ln4\pi.
\end{equation}
Keep \(\tr\mathbf1=4\) as the Dirac trace convention while continuing
contracted Lorentz indices to \(d\). In particular,
\(\gamma^\mu\gamma_\mu=d\) and
\(\gamma^\mu\slashed a\gamma_\mu=(2-d)\slashed a\).
Terms proportional to \(\eps\) can multiply \(1/\eps\) poles and give
finite contributions. They must not be discarded prematurely.

Introduce \(A_0=\sqrt{Z_3}A\), \(\psi_0=\sqrt{Z_2}\psi\),
\(m_0=Z_m m\), and
\(e_0=\mu^\eps eZ_1/(Z_2\sqrt{Z_3})\). Writing
\(Z_i=1+\delta_i\), \(Z_m=1+\delta Z_m\), the one-loop counterterms are
\begin{equation}
 \Lag_{\rm ct}=-\frac{\delta_3}{4}F^2
 +\delta_2\bar\psi(\ii\slashed\partial-m)\psi
 -m\delta Z_m\bar\psi\psi
 -\mu^\eps e\delta_1\bar\psi\gamma^\mu\psi A_\mu.
 \label{eq:ct}
\end{equation}
Products of counterterms enter only at higher orders. Choosing
\(\xi_0=Z_3\xi\) leaves the displayed gauge-fixing term in renormalized
form. A gauge-preserving subtraction will imply \(Z_1=Z_2\).

The loop program is now concrete. First regulate and evaluate a 1PI
function. Next cancel its local ultraviolet pole using \eqref{eq:ct}.
Finally impose a subtraction condition, such as \(\MSbar\) or a physical
on-shell condition, and combine the result with the other contributions
to an observable. A subtraction prescription changes parameter labels,
not predictions expressed in the same measured inputs.

\begin{remarkbox}[Ultraviolet and infrared limits answer different questions]
A divergence from large loop momentum fixes local counterterms. A
divergence from a photon with vanishing energy reflects the impossibility
of resolving arbitrarily soft radiation. The latter is handled in
inclusive observables or suitably dressed charged states, not by a new
ultraviolet counterterm. Keep the two regulators conceptually distinct,
even when dimensional regularization represents both by poles.
\end{remarkbox}
Lecture 28 computes the propagator corrections and their physical
interpretation. Lecture 29 completes the vertex, explains why the same
charge appears in propagation and interaction, and extracts the finite
anomalous magnetic moment. Homework 9 ties these calculations together.
\end{document}
